\chapter{Understanding Data Institutions}

Bottom-up DIs are a specific form of commons-based data governance, often predicated on a fiduciary or delegated relationship between beneficiaries and their representatives (\cite{delacroix2020DemocratisingDigitalRevolution}). As a parsimonious institutional template, these DIs offer an extensible model of democratic data governance. I consider DI use cases narrowly applied to mitigating existing harms and generating public value throughout the stages of ML dataset development.


The crux of this thesis (covered next section) seeks to extend the set of proposed DI use cases beyond those already well-theorised by the institutional literature's emphasis on DIs' \textit{upstream} institutional logic (covered this section). This upstream logic primarily revolves around redistributing value across imbalanced digital markets. Additionally, I propose that DIs offer further dimensions of ML-specific utility, which become apparent when incorporating technical and policy-based critiques of current ML dataset development stemming from other fields of study.

Current redistribution-focused DI use cases include: reappropriating value extracted from disempowered data producers (\cite{birkinbine2018CommonsPraxisCritical, loi2023RawlsianPropertyowningDemocracy, spiekermann2021BigDataJustice}), opening up new pockets of value across the digital economy to encourage greater data sharing (\cite{2023UnderstandingSocialEconomic}), structuring a digital commons (\cite{delacroix2019BottomupDataTrusts, singh2019DataDigitalIntelligence}), or supporting data availability for impactful value-aligned projects (\cite{baarbe2019ProposedAgriculturalData}). These remain important aspects of DIs' utility to the ML community, but I further propose that establishing dedicated DIs for the design, development and stewarding of ML datasets creates integrative opportunities for both upstream \textit{and downstream} value. To make this case, in the next section I draw on knowledge and critiques generated by the data management and (critical) global governance literatures to additionally emphasise DIs' value proposition to the ML community, impacted stakeholders (including data subjects) and wider society. But first, a genealogy of DI theorising to illustrate that just as ML datasets represent strategic policy objects, so DIs represent a strategic institutional template for dataset stewarding.

\section*{Intellectual foundations}
\addcontentsline{toc}{section}{Intellectual foundations}

DI theorising starts with recognising the uniqueness of data as a resource, conceptualising it as a non-rival, relational public/club good, replete with informational externalities and economies of scale (in collection and inference). This lends data well to \textit{commons}-style governance arrangements.\footnote{Theorising around a ``digital commons" represents a major governance topic in and of itself. I consider it here only to the extent that it is relevant and adjacent to discussion of DIs. The following section explicitly considers ML's relationship to the commons, but for more, see \textcite{dulongderosnay2020DigitalCommons, huang2023GenerativeAIDigital, ostromGoverningCommonsEvolution, dietz2003StruggleGovernCommonsa, grossman2016CaseDataCommons, baarbe2019ProposedAgriculturalData, ottolia2022TopographyDataCommons, spiekermann2021BigDataJustice, meadwayCREATINGDIGITALCOMMONS, dulongderosnay2020DigitalCommons, potts2023ProfitingDataCommons, yakowitzTRAGEDYDATACOMMONS})} These unique properties of data serve as the theoretical foundations for developing and evaluating the institutional and legal arrangements of digital governance efforts (\cite{mills2019WhoOwnsFuture}). 

DI theorising was for a while subsumed under the label of data trusts, now that relationship is reversed. Yet it remains instructive to consider the development and diversity of input into DI theorising by tracing the development of its flagship institutional form. As the most common private law instrument for fiduciary management of assets in the public interest, trusts present data theorists with a tenable instrument for post-hoc insertion of public value considerations into the governance of increasingly privatised digital infrastructures (\cite{mcdonaldReclaimingDataTrusts}). Most commonly, data trusts are now proposed as a response to the growing alarm around harms deriving from the emerging `platform model' (see \textcite{birkinbine2018CommonsPraxisCritical, meadwayCREATINGDIGITALCOMMONS, coyle2021PotentialSocialValuea}), although this was not always the case.

As a legal instrument structured around the appointment of a steward/trustee to manage data on beneficiaries’ behalf (or the rights to data) for a predefined purpose (\cite{ohara2020DataTrusts}), data trusts have periodically been proposed as the solution to various imbalances associated with digital markets. But trusts were not always presented as a means for representing the public interest, as they often are today. Early DI proponents innovated trusts for brokering individual digital privacy risks (\cite{edwards2004ProblemPrivacyModest}) and later as (libertarian-esque) platform fiduciaries (\cite{BalkinFiduciary}). Ironically, recent critiques of individualist notions of data rights (see \textcite{viljoen2021RelationalTheoryData}) have also led to calls for data trusts in open data projects, whereby trusts are touted as a governance arrangement able to incentivise greater sharing and public interest leveraging of data’s aggregate utility (\cite{2021EconomicImpactTrust}). \textcite{delacroix2019BottomupDataTrusts} persuasively developed data-trusts-as-legal-vehicle theorising, calling for a trust ecosystem constituted of differentiated units, each catering to a specific subset of governance demands and customised to community values. At a systems-level, they propose this market for trusts could accommodate a range of heterogeneous data contexts. It is this `open data' line of reasoning that now informs DI theorising more broadly, as DIs are put forward as a form of institutional redress to imbalanced capitalistic practice undermining the digital commons.

Proponents of open data propose DIs as the structuring unit upon which to collectively steward an informational commons accountable to data-generating communities and respecting their rights and interests. These proponents range from business-oriented \textcite{zarkadakis2020DataTrustsCould}, to privacy activist organisations (\cite{DataTrustsInitiative}), to Global South stakeholders (\textcite{borokini2021ExploringthefutureofdatagovernaceinAfricaDataStewardship}), and social enterprises (see \cite{CasesDataCollaboratives}). As the domain has matured, a plethora of alternative names and overlapping institutional propositions beyond trusts have emerged, yet the key insight of delegated fiduciary responsibility as a means to leverage data’s aggregative properties, remains core.

\section*{Contested frames}
\addcontentsline{toc}{section}{Contested frames}

DIs occupy a contested space. An early proponent of public-interest DIs was the UK government, which identified data trusts as a suitable governance instrument for data sharing infrastructure in an AI economy (\cite{hall2017GrowingArtificialIntelligence}). Responding in part to the backlash against the sharing of NHS data with DeepMind, the 2017 report presented an early emphasis on downstream value creation that has since waned. Since then, the public dialogue is increasingly shaped by well-financed interests looking to maximise large-scale privately-driven data sharing (as motivated the controversy around Alphabet's Sidewalk Labs - now closed). In this framing, data trusts’ privacy-enhancing functions are variously imagined as atomising, anarcho-libertarian alternatives to state control of data resources, or as disingenuous big tech cover for the freer flow and private capture of information (see \textcite{bietti2021GenealogyDigitalPlatform}). 

More recently, open data proponents have taken up the DI mantle, calling for the public sector to catch up and engage (\cite{RoleGovernmentProvider, 2022HowCanUK}). The point of departure for my proposed use of DIs for ML datasets are the variants of this discourse that frames DIs as solving the \textit{paradox of openness} posed by extractive platforms (\cite{jameson2019ExtractiveLogicsSystemic, verdegem2022DismantlingAICapitalism}), and providing \textit{the} structuring institution for the digital commons as a shared intellectual inheritance (\cite{khong2023ARTIFICIALINTELLIGENCECOMMON}). In this `open data' framing, a sustainable AI economy requires addressing the associated harms, unequal distributions of gains, and concerns about future precarity. DIs offer one such institutional means to do so; as the foundation for AI governance, inclusive data governance will increasingly form a bedrock social contract (\cite{InterwovenRealmsData}). 

Beyond Tarkowski and Warso's March 2024 contribution, there have been few existing initiatives to articulate the use of collective data institutions for ML datasets, although several works allude to the possibility (see \textcite{delacroix2020DemocratisingDigitalRevolution, huang2023GenerativeAIDigital, zygmuntowski2021EmbeddingEuropeanValues}). The closest is \textcite{chan2023ReclaimingDigitalCommons}, but their focus is \textit{national}, while \cite{jernite2022DataGovernanceAge} propose a data stewardship organisation to establish and formalise relationships between actors in data ecosystems at the \textit{international} level. On the contrary, my focus is the \textit{community} level.

